\documentclass[a4paper,9pt]{extarticle}

%=========================================================
% PACKAGES
%=========================================================

\usepackage[
a4paper,
top=1.8cm,
bottom=1.5cm,
footskip=7mm,
left=1cm,
right=1cm,
headheight=14pt,
headsep=5mm
]{geometry}
\usepackage[T1]{fontenc}
\usepackage[utf8]{inputenc}
\usepackage[french]{babel}

\usepackage[sfdefault]{FiraSans}
\renewcommand{\familydefault}{\sfdefault}

\usepackage{multicol}
\setlength{\columnsep}{0.7cm}

\usepackage{amsmath,amssymb}

\usepackage{titlesec}
\usepackage{enumitem}

\usepackage{tcolorbox}
\tcbuselibrary{skins,breakable}

\usepackage{xcolor}

\usepackage{graphicx}
\usepackage{float}

\usepackage{tikz}

\usepackage{fancyhdr}
\usepackage[hidelinks]{hyperref}


%=========================================================
% PAGE SETTINGS
%=========================================================

\setlength{\parindent}{0pt}
\setlength{\parskip}{1.5pt}

\setlength{\abovedisplayskip}{4pt}
\setlength{\belowdisplayskip}{4pt}

\setlist[itemize]{
leftmargin=*,
itemsep=1pt,
topsep=1pt
}


%=========================================================
% COLORS
%=========================================================

\definecolor{Blue}{HTML}{E3F0FF}
\definecolor{Green}{HTML}{E3F7EA}
\definecolor{Pink}{HTML}{FFE5ED}
\definecolor{Orange}{HTML}{FFF0DC}
\definecolor{Purple}{HTML}{EEE7FF}
\definecolor{Yellow}{HTML}{FFFADB}
\definecolor{Teal}{HTML}{E0F5F3}


\definecolor{BlueTitle}{HTML}{285B8F}
\definecolor{GreenTitle}{HTML}{287052}
\definecolor{PinkTitle}{HTML}{A84063}
\definecolor{OrangeTitle}{HTML}{A96500}
\definecolor{PurpleTitle}{HTML}{6245A5}
\definecolor{TealTitle}{HTML}{1F6F68}


%=========================================================
% HEADER
%=========================================================

\pagestyle{fancy}

\fancyhf{}

\fancyhead[L]{\small\textbf{Modèles stochastiques}}
\fancyhead[C]{\small Wiener, prédiction linéaire et Poisson}
\fancyhead[R]{\small Week 09 --- 02-08/11/2026}

\fancyfoot[C]{\small\href{https://github.com/AlphaXT38G/notes_ba5}{github.com/AlphaXT38G/notes\_ba5}}
\fancyfoot[R]{\small \thepage}

\renewcommand{\headrulewidth}{0.4pt}



%=========================================================
% TITLES
%=========================================================

\titleformat{\section}
{\large\bfseries\color{BlueTitle}}
{}{0pt}{}

\titlespacing{\section}
{0pt}{6pt}{3pt}


% separator

\newcommand{\separator}{
\vspace{2pt}
\hrule
\vspace{3pt}
}


%=========================================================
% BOXES
%=========================================================

% Boxes are kept whole (unbreakable): inside multicols, a box split by tcolorbox
% leaves an empty column behind it.

% DEFINITION - blue
\newtcolorbox{definition}{
enhanced,
unbreakable,
colback=Blue,
colframe=BlueTitle,
boxrule=0.7pt,
arc=0mm,
title=\textbf{Definition},
fonttitle=\bfseries
}

% IMPORTANT - pink
\newtcolorbox{important}{
enhanced,
unbreakable,
colback=Pink,
colframe=PinkTitle,
boxrule=0.7pt,
arc=0mm,
title=\textbf{Key idea},
fonttitle=\bfseries
}

% FORMULAS - green
\newtcolorbox{formula}{
enhanced,
unbreakable,
colback=Green,
colframe=GreenTitle,
boxrule=0.7pt,
arc=0mm,
title=\textbf{Key formulas},
fonttitle=\bfseries
}

% WARNING - orange
\newtcolorbox{warning}{
enhanced,
unbreakable,
colback=Orange,
colframe=OrangeTitle,
boxrule=0.7pt,
arc=0mm,
title=\textbf{Warning},
fonttitle=\bfseries
}

% EXAMPLE - purple
\newtcolorbox{example}{
enhanced,
unbreakable,
colback=Purple,
colframe=PurpleTitle,
boxrule=0.7pt,
arc=0mm,
title=\textbf{Example},
fonttitle=\bfseries
}

% SUMMARY - orange
\newtcolorbox{summary}{
enhanced,
unbreakable,
colback=Yellow,
colframe=OrangeTitle,
boxrule=0.7pt,
arc=0mm,
title=\textbf{Summary},
fonttitle=\bfseries
}



% SIMPLE - teal: the same idea again, in the plainest words, with an everyday picture
\newtcolorbox{simple}{
enhanced,
unbreakable,
colback=Teal,
colframe=TealTitle,
boxrule=0.7pt,
arc=0mm,
title=\textbf{Expliqué simplement},
fonttitle=\bfseries
}



%=========================================================
% FIGURE BOX
%=========================================================

\newtcolorbox{figurebox}{
enhanced,
colback=white,
colframe=gray,
boxrule=0.5pt,
arc=0mm,
center
}



%=========================================================
% DOCUMENT
%=========================================================

\begin{document}

%\footnotesize   --> To reduce the size of the text

\thispagestyle{fancy}



\begin{multicols}{2}


%=========================================================
\section{Où on en est}


\begin{summary}

\begin{itemize}

\item Semaine 7 : break. Semaine 8 : office hours le mardi, \textbf{midterm} (modules 1 à 3)
le vendredi. Pas de nouvelle matière, donc pas de notes pour ces deux semaines.

\item Mardi : module 4 (fin), filtre de Wiener et prédiction linéaire ; module 5, processus de
Poisson : définition, équations de Kolmogorov, loi.

\item Vendredi : module 5, propriétés, bruit impulsif de Poisson (\emph{shot noise}).

\item Exercices M4 : 1, 4 (mardi) ; 5, 8, 9 (vendredi).

\end{itemize}

\end{summary}


\begin{important}

\textbf{Conseils du Moodle pour l'exercice M4 4.}

\begin{itemize}

\item \textbf{4 b} : s'arrêter à une expression avec une \textbf{double somme}. La réduire à la
somme simple de l'énoncé (comme dans les solutions) est technique et hors exercice.

\item \textbf{4 c} : partir des formules données plus haut dans l'énoncé.

\end{itemize}

\end{important}


%=========================================================
\section{Rappel : le filtre de Wiener}


\begin{formula}

Entrée $X(n)$ WSS, sortie $Y(n) = \sum_kh(k)X(n-k)$, réponse désirée $D(n)$, erreur $e(n) =
D(n) - Y(n)$, on minimise $\varepsilon = E[e^2(n)]$ :

\[
\sum_lh(l)R_X(k-l) = R_{DX}(k) \quad\Rightarrow\quad H(f) = \frac{S_{DX}(f)}{S_X(f)}
\]

(non causal), ou, en causal, les \textbf{équations de Wiener-Hopf} : $\sum_{l\ge0}h(l)R_X(k-l)
= R_{DX}(k)$ pour $k\ge0$.

\end{formula}


\begin{important}

\textbf{Exercice 8} (filtre de prédiction ci-dessous) : l'erreur $e(n)$ et l'entrée $X$ sont
\textbf{orthogonales}. C'est ce que disent les équations : pour chaque $k$ où le filtre a un
coefficient, $E[e(n)X(n-k)] = R_{DX}(k) - \sum_lh(l)R_X(k-l) = 0$.

\end{important}


\begin{simple}

Le meilleur filtre est celui dont l'erreur restante n'a \textbf{plus rien en commun} avec ce
qu'on a vu à l'entrée : s'il restait un lien, on pourrait s'en servir pour corriger encore un
peu l'erreur, et le filtre ne serait pas le meilleur.

\end{simple}


%=========================================================
\section{Application : restauration d'image}


\begin{definition}

Modèle de dégradation (« image » à une dimension pour commencer) :

\[
X(n) = \sum_ib(i)S(n-i) + N(n)
\]

$S$ : image idéale, $X$ : image dégradée, $b$ : filtre de dégradation (mouvement de la
caméra...) de transmittance $B(f)$, $N$ : bruit additif. On cherche $H(f)$ tel que la sortie
$Y$ soit proche de $S$.

\end{definition}


\begin{warning}

\textbf{Filtrage inverse} $H_{\mathrm{inv}}(f) = 1/B(f)$ : $\tilde Y(f) = \tilde S(f) +
\tilde N(f)/B(f)$ (transformées de Fourier des signaux, pas des densités spectrales). Si $B(f)$
est très petit à certaines fréquences, le bruit $\tilde N/B$ explose et écrase l'image. Et si
l'image est seulement bruitée ($B = 1$), il ne fait rien.

\end{warning}


\begin{formula}

\textbf{Wiener}, avec $D = S$, $N$ centré et indépendant de $S$ :

\[
S_{DX}(f) = B^\star(f)S_S(f) \qquad S_X(f) = |B(f)|^2S_S(f) + S_N(f)
\]
\[
H(f) = \frac1{B(f)}\cdot\frac{|B(f)|^2S_S(f)}{|B(f)|^2S_S(f) + S_N(f)}
\]

Une cascade : un filtre qui \textbf{réduit le bruit}, puis le \textbf{filtre inverse}. Sans
bruit, c'est le filtre inverse ; avec $B = 1$, c'est le Wiener qui débruite. En 2D, même formule
en $(f_1,f_2)$.

\end{formula}


\begin{important}

Limites : le critère des moindres carrés n'est pas celui de l'\oe il (variantes avec des
paramètres $\alpha$, $\beta$), une image n'est pas stationnaire (version adaptative, algorithme
LMS), et il faut connaître $S_S$ et $S_N$.

\end{important}


\begin{simple}

Une photo floue et granuleuse. Le filtre inverse « défloute » à fond, mais il amplifie aussi
le grain là où le flou avait presque tout effacé : catastrophe. Wiener est plus malin : il ne
défloute que là où le signal est plus fort que le bruit, et laisse tomber les fréquences où il
n'y a plus que du grain.

\end{simple}


%=========================================================
\section{Application : codage PCM et DPCM}


\begin{definition}

\textbf{PCM} (MIC) : on quantifie chaque échantillon directement, en l'arrondissant au $i\Delta$
le plus proche ($\Delta$ : pas de quantification). Simple et rapide, mais n'exploite pas la
redondance : environ 64 kbps pour la parole.

\vspace{3pt}

\textbf{Bruit de quantification} $q(n) = X(n) - Q(X(n))$ (en négligeant les dépassements de
plage) : sous des hypothèses raisonnables, uniforme sur $[-\Delta/2,\Delta/2]$, indépendant
d'un échantillon à l'autre, donc un \textbf{bruit blanc} :

\[
\sigma_q^2 = \frac{\Delta^2}{12} \qquad R_q(k) = \sigma_q^2\delta(k) \qquad \mathrm{SNR} =
\frac{\sigma_X^2}{\sigma_q^2}
\]

\end{definition}


\begin{definition}

\textbf{DPCM} : on \textbf{prédit} $X(n)$ à partir des échantillons précédents, $Y(n)$, et on ne
quantifie et transmet que l'erreur de prédiction $e(n) = X(n) - Y(n)$. Le récepteur refait la
même prédiction et ajoute l'erreur reçue.

\vspace{3pt}

Le prédicteur minimise la variance de $e$, donc $\sigma_e^2 \ll \sigma_X^2$ : pour le même
nombre de bits, bruit de quantification bien plus faible ; pour la même qualité, moins de bits
(environ 16 kbps).

\end{definition}


\begin{important}

En pratique, on prédit à partir des échantillons \textbf{quantifiés} $X_q(n-1),\dots$ (le
prédicteur est dans la boucle autour du quantificateur) : les erreurs de quantification sont
prises en compte et ne s'accumulent pas.

\end{important}


\begin{simple}

Au lieu d'envoyer chaque valeur en entier, on envoie seulement « \textbf{la surprise} » : la
différence entre la vraie valeur et ce qu'on pouvait deviner. Comme la parole change
doucement, la surprise est petite, donc il faut moins de bits pour l'écrire.

\end{simple}


%=========================================================
\section{Filtre de prédiction linéaire}


\begin{formula}

$X(n)$ WSS. On estime $X(n)$ à partir des $m$ échantillons précédents :

\[
Y(n) = \sum_{k=1}^mh(k)X(n-k)
\]

en minimisant $\varepsilon = E[(X(n) - Y(n))^2]$ : Wiener avec $D(n) = X(n)$, donc $R_{DX}(k)
= R_X(k)$, filtre à réponse impulsionnelle finie. Les équations de Wiener-Hopf deviennent les
\textbf{équations de Yule-Walker} :

\[
\sum_{l=1}^mh(l)R_X(k-l) = R_X(k), \quad 1\le k\le m
\]

Avec $\mathbf r_X = [R_X(1)\dots R_X(m)]^T$ et $\mathbf h = [h(1)\dots h(m)]^T$ :

\[
\Phi_X\mathbf h = \mathbf r_X \qquad \mathbf h = \Phi_X^{-1}\mathbf r_X
\]

\end{formula}


\begin{important}

Il faut inverser la matrice de corrélation $\Phi_X$ (symétrique, Toeplitz).

\begin{itemize}

\item Symétrique : $\Phi_X = Q^T\Lambda Q$ avec $Q$ orthogonale (décomposition de
\textbf{Karhunen-Loève}, qui décorrèle les échantillons), $\Phi_X^{-1} = \sum_i\lambda_i^{-1}
\mathbf q_i\mathbf q_i^T$. Élégant, mais calculer les valeurs propres coûte cher.

\item En pratique : l'algorithme de \textbf{Levinson}, qui exploite la structure de Toeplitz.

\end{itemize}

\end{important}


\begin{simple}

Deviner la prochaine valeur comme un mélange des $m$ dernières. Les bons poids $h(k)$ se
trouvent en résolvant un petit système linéaire, qui ne demande que l'auto-corrélation : à
quel point une valeur ressemble à celle d'il y a 1, 2, ..., $m$ pas.

\end{simple}


%=========================================================
\section{Application : vocodeur LPC}


\begin{definition}

\textbf{Synthèse de la parole} (modèle simplifié) : un filtre AR de coefficients $a_k$, excité
par $u(n)$ :

\begin{itemize}

\item sons \textbf{voisés} (/a/, /i/, /v/) : train périodique d'impulsions d'amplitude $\sigma$
(les cordes vocales), de période le \textbf{pitch} ;

\item sons \textbf{non voisés} (/f/, /s/, /t/) : bruit blanc de variance $\sigma^2$.

\end{itemize}

Sur des blocs courts (10 à 30 ms), la parole est considérée stationnaire.

\end{definition}


\begin{formula}

\textbf{Vocodeur} (\emph{voice coder}) : pour chaque bloc, l'émetteur ne transmet que : voisé
ou non, le pitch, le gain $\sigma$, et les coefficients de prédiction $h(k)$ (estimées des
$-a_k$, ordre $m = p$). Le récepteur resynthétise la parole avec le modèle.

\vspace{3pt}

Débit d'environ 4.8 kbps, voire 2.4 kbps (au lieu de 16), mais la parole devient
\textbf{synthétique}.

\end{formula}


\begin{simple}

Au lieu d'envoyer le son, on envoie la \textbf{recette} pour le refaire : « une voix grave, qui
dit un son en /a/, à tel volume ». À l'arrivée, une petite machine reproduit le son à partir
de la recette. Très économique, mais ça sonne comme un robot.

\end{simple}


%=========================================================
\section{Processus de comptage}


\begin{definition}

$\{N(t), t\in\mathbb R^+\}$ est un \textbf{processus de comptage} si $N(t)$ est le nombre
d'évènements dans $[0,t]$ :

\begin{enumerate}[leftmargin=*,itemsep=1pt,topsep=1pt]

\item temps continu, valeurs entières $\ge0$ ;

\item $N(0) = 0$ ;

\item $t_1 < t_2 \Rightarrow N(t_1) \le N(t_2)$ ;

\item le nombre d'évènements dans $]t_1,t_2]$ est $N(t_2) - N(t_1)$.

\end{enumerate}

\end{definition}


\begin{formula}

Deux processus associés (temps discret, valeurs continues) :

\begin{itemize}

\item $S(n)$ : \textbf{temps d'arrivée} du $n$-ième évènement ;

\item $T(n)$ : \textbf{temps entre arrivées}, entre le $n$-ième et le $(n+1)$-ième.

\end{itemize}

\[
N(t)\ge n \Leftrightarrow S(n)\le t \qquad S(n) = \sum_{m=0}^{n-1}T(m)
\]
\[
T(n) = S(n+1) - S(n) \;(n\ge1),\quad T(0) = S(1)
\]
\[
N(t) = \max\{n : S(n)\le t\} \qquad S(n) = \inf\{t : N(t) = n\}
\]

\end{formula}


\begin{simple}

Un compteur à l'entrée d'un magasin. $N(t)$ : combien de clients sont entrés depuis
l'ouverture. $S(n)$ : à quelle heure est entré le $n$-ième. $T(n)$ : combien de temps on a
attendu entre deux clients. Trois façons de raconter la même histoire.

\end{simple}


%=========================================================
\section{Processus de Poisson : deux définitions}


\begin{definition}

\textbf{Définition 1.} Poisson homogène d'intensité $\lambda > 0$ : un processus de comptage
avec

\begin{itemize}

\item \textbf{H1.} \textbf{accroissements indépendants} (des intervalles disjoints sont
indépendants) ;

\item \textbf{H2.} \textbf{homogène} (accroissements stationnaires) : $P(N(t+T) - N(t) = n_0) =
P(N(T) = n_0)$ ;

\item \textbf{H3.} pas d'arrivées groupées :

\[
P(N(\Delta t) = 1) = \lambda\Delta t + o(\Delta t)
\]
\[
P(N(\Delta t)\ge2) = o(\Delta t)
\]

donc $P(N(\Delta t) = 0) = 1 - \lambda\Delta t + o(\Delta t)$ ($f$ est $o(\Delta t)$ si
$f(\Delta t)/\Delta t\to0$).

\end{itemize}

\end{definition}


\begin{definition}

\textbf{Définition 2.} H1, et \textbf{H2'} : pour tous $t, T\ge0$,

\[
P(N(t+T) - N(t) = n) = \frac{(\lambda T)^n}{n!}e^{-\lambda T}
\]

\textbf{Théorème 1} : les deux définitions sont équivalentes.

\end{definition}


\begin{simple}

Des clients qui arrivent \textbf{complètement au hasard} : ce qui se passe dans un intervalle
n'influence pas un autre (H1), toutes les heures se valent (H2), et deux clients n'arrivent
jamais pile en même temps (H3). Ces trois règles de bon sens suffisent à imposer que le nombre
de clients dans une durée $T$ suit une loi de Poisson de moyenne $\lambda T$.

\end{simple}

\end{multicols}


\end{document}
